\chapter[INTRODUCTION À APRS]{Introduction à APRS}\label{introduction-uxe0-aprs}

\phantomsection\addcontentsline{toc}{section}{Qu'est-ce qu'APRS ?}
\begin{aprssideblock}{Qu'est-ce\\qu'APRS ?}
APRS est l'abréviation d'\emph{Automatic Position Reporting System},
conçu par Bob Bruninga WB4APR et présenté par lui à la conférence
TAPR/ARRL Digital Communications de 1992.

Fondamentalement, APRS est un protocole de communication packet destiné
à diffuser en temps réel des données vivantes à tous les nœuds d'un
réseau. Sa fonctionnalité la plus visible est la combinaison du packet
radio avec le réseau satellite GPS, permettant aux radioamateurs
d'afficher automatiquement sur une carte PC la position de stations et
d'autres objets. D'autres fonctions non directement liées au reporting
de position sont supportées : stations météo, radiogoniométrie et
messagerie.

APRS diffère du packet classique à plusieurs égards :

\begin{itemize}
\tightlist
\item
  Cartes et affichages de données, pour la localisation
  véhicules/personnels et le reporting météo en temps réel.
\item
  Communications \emph{one-to-many}, chacun étant mis à jour
  immédiatement.
\item
  Digipeating générique avec alias d'indicatifs bien connus --- pas
  besoin de connaître la topologie du réseau à l'avance.
\item
  Digipeating intelligent avec substitution d'indicatif pour réduire le
  flooding.
\item
  Sur trames AX.25 UI, supporte la messagerie bidirectionnelle et la
  diffusion de bulletins/annonces.
\item
  Communications avec les radios Kenwood TH-D7 et TM-D700 (TNC et
  firmware APRS intégrés).
\end{itemize}

Le packet radio conventionnel n'est vraiment utile que pour transporter
du trafic messagerie en masse d'un point à un autre, et a
traditionnellement été difficile à appliquer à des événements temps réel
où l'information a une durée de vie très courte. APRS transforme le
packet radio en un système de communication et d'affichage tactique
temps réel pour urgences et missions de service public.

APRS fournit une connectivité universelle entre toutes les stations,
sans la complexité, les délais et les limites d'un réseau à connexion.
N'importe quel nombre de stations peut échanger des données comme le
feraient des utilisateurs voix sur un net vocal. Toute station qui a
quelque chose à partager l'envoie, toutes les stations le reçoivent et
le journalisent.

APRS reconnaît qu'un des besoins temps réel les plus importants lors
d'un événement ou d'une urgence est le suivi d'actifs clés. Où est le
leader du marathon ? Où sont les véhicules de secours ? Quelle météo en
tel point ? Où sont les lignes électriques tombées ? Où est la tête du
défilé ? Où est la caméra ATV mobile ? Où est l'orage ?

Pour ces questions, APRS fournit un système complet de localisation
automatique et de reporting de statut. Il peut fonctionner sur n'importe
quel système radio bidirectionnel --- amateur, marine, cellulaire. Il
existe même un réseau international de suivi APRS en direct sur
Internet.
\end{aprssideblock}

\phantomsection\addcontentsline{toc}{section}{Fonctionnalités APRS}
\begin{aprssideblock}{Fonctionnalités\\APRS}
APRS tourne sur la plupart des plateformes : DOS, Windows 3.x, Windows
95/98, MacOS, Linux, Palm. La plupart des implémentations supportent les
principales fonctions :

\begin{itemize}
\tightlist
\item
  \textbf{Cartes} --- positions des stations APRS tracées en temps réel,
  échelle de quelques centaines de mètres à monde entier. Les stations
  reportant course et vitesse sont extrapolées à leur position actuelle
  par dead-reckoning. Bases de données superposables : digipeaters APRS,
  sites du National Weather Service US, magasins de matériel radio. Zoom
  possible sur n'importe quel point du globe.
\item
  \textbf{Reporting de station météo} --- affichage automatique
  d'informations météo distantes à l'écran.
\item
  \textbf{Reporting DX Cluster} --- outil idéal pour l'utilisateur de DX
  cluster. Quelques stations APRS connectées à des clusters peuvent
  relayer les infos de station DX à beaucoup d'autres stations locales,
  réduisant la charge packet globale sur les clusters.
\item
  \textbf{Accès Internet} --- Internet peut être utilisé de manière
  transparente pour interconnecter des réseaux radio locaux partout dans
  le monde. On peut se connecter en telnet à des serveurs APRS Internet
  et voir en direct des centaines de stations du monde entier. Tout le
  monde peut alimenter le serveur avec ses paquets locaux, tout le monde
  partout peut les voir.
\item
  \textbf{Messages} --- messages bidirectionnels avec accusé de
  réception. Chaque message entrant alerte l'utilisateur et reste
  affiché jusqu'à effacement manuel.
\item
  \textbf{Bulletins et annonces} --- adressés à tous. Les bulletins sont
  envoyés quelques fois par heure pendant quelques heures ; les annonces
  moins fréquemment mais peut-être sur plusieurs jours.
\item
  \textbf{Suivi de stations fixes} --- en plus du suivi automatique des
  stations mobiles GPS/LORAN, APRS suit aussi les rapports manuels ou
  par grid square.
\item
  \textbf{Objets} --- n'importe quel utilisateur peut placer un Object
  APRS sur sa carte, et en quelques secondes cet objet apparaît sur tous
  les autres écrans. Utile pour tracker des ressources non équipées de
  trackers : un seul opérateur, en écoutant la voix, tient à jour les
  positions ; toutes les autres stations en profitent.
\end{itemize}
\end{aprssideblock}

